\section{evnx auth}

\begin{frame}{\cmd{evnx auth} — your evnx cloud account}
  \begin{block}{Synopsis}\texttt{evnx auth <SUBCOMMAND>}\end{block}

  \begin{tabular}{@{}ll@{}}
    \toprule
    \texttt{auth register} & Create an account (keys generated locally) \\
    \texttt{auth login}    & SRP-6a login, TOTP challenge if enabled \\
    \texttt{auth logout}   & End the session \\
    \texttt{auth status}   & Who am I, against which server \\
    \midrule
    \texttt{auth token create/list/revoke} & CI tokens (\texttt{evnx\_tok\_\ldots}) \\
    \texttt{auth totp enable/disable}      & Second factor \\
    \texttt{auth totp recovery-codes}      & 10 single-use codes, shown once \\
    \texttt{auth sessions list/revoke}     & Active sessions \\
    \texttt{auth sessions revoke-others}   & After a laptop goes missing \\
    \bottomrule
  \end{tabular}

  \vspace{0.5em}
  {\footnotesize Credentials live in \file{\textasciitilde/.config/evnx/},
  mode \texttt{0600} — the \file{\textasciitilde/.aws/credentials} model.}
\end{frame}

\begin{frame}[fragile]{The login flow — your password never leaves the device}
\begin{semiverbatim}\scriptsize
POST /auth/register    SRP verifier + salts + public keys (all client-side)
GET  /auth/verify-email   emailed link flips email_verified
POST /auth/srp/init    -> session_id, salts, server public B
POST /auth/srp/verify  M1 -> M2 + \{ access, refresh \}
                          -> M2 + totp_pending_token   (if TOTP on)
POST /auth/totp/verify pending token + code -> \{ access, refresh \}
\end{semiverbatim}

  \vspace{0.4em}
  \begin{block}{SRP-6a in one sentence}
    The server stores a \emph{verifier}, not a password hash it could brute
    force offline in the usual way, and the proof exchange never transmits the
    password or anything derived from it that a passive observer could reuse.
  \end{block}

  {\footnotesize 15-minute access token; 30-day refresh token, rotated on use.}
\end{frame}

\begin{frame}[fragile]{The transport guard}
  \capture{66-tls-guard.txt}

  {\footnotesize Loopback is exempt, so local development against
  \texttt{http://127.0.0.1:8080} works. The error explains the \emph{threat},
  not just the rule.}
\end{frame}

\begin{frame}{Use cases}
  \begin{itemize}\footnotesize
    \item \textbf{CI token, scoped to one vault, read-only}\\
      \texttt{evnx auth token create --scope read\_only --vault app/production}
    \item \textbf{Turn on 2FA} — \texttt{evnx auth totp enable} (QR in the terminal)
    \item \textbf{Save the recovery codes} — \texttt{evnx auth totp recovery-codes}
    \item \textbf{Lost laptop} — \texttt{evnx auth sessions revoke-others}
  \end{itemize}

  \vspace{0.8em}
  \begin{block}{Minting a token requires a real login}
    \texttt{/auth/tokens} sits behind a user-session guard, so a leaked CI token
    cannot mint itself a longer-lived, wider-scoped replacement and survive
    revocation of the original.
  \end{block}
\end{frame}
